\clearpage
\VerdierPhysicalPageAnchor{0066}%
\VerdierSemanticAnchor{ch1-2}
\section{Complexes d'une catégorie additive}

\VerdierSemanticAnchor{ch1-2-1}
\subsection{Complexes \texorpdfstring{$n$}{n}-uples}

\VerdierSemanticAnchor{ch1-2-1-1}
\noindent\textbf{2.1.1.}\quad
Fixons d'abord quelques notations. On désigne par $[n]$ l'ensemble :

\VerdierSemanticAnchor{ch1-2-1-1-set}
\begin{equation*}
  [n]=\{p\in\mathbb Z\mid 1\leq p\leq n\}.
\end{equation*}

\VerdierSemanticAnchor{ch1-2-1-1-free-group}
On désigne par $\mathbb Z[n]$ le groupe abélien libre engendré par
l'ensemble $[n]$ et pour tout $i\in[n]$ par $e_i$ le générateur de
$\mathbb Z[n]$ associé à $i$. Soit $\varphi\colon[n]\longrightarrow[m]$
une application ; on désigne encore par
$\varphi\colon\mathbb Z[n]\longrightarrow\mathbb Z[m]$ l'homomorphisme
de groupes qui lui correspond. Dans ce paragraphe, $C$ désigne une
catégorie additive.

\VerdierSemanticAnchor{ch1-2-1-2}
\VerdierSemanticAnchor{ch1-2-1-2-definition}
\noindent\textbf{Définition 2.1.2.}\quad
Un \emph{complexe $n$-uple d'objets de $C$} est un objet constitué par :

\begin{enumerate}
\renewcommand{\labelenumi}{\arabic{enumi})}
\item \VerdierSemanticAnchor{ch1-2-1-2-objects}%
  La donnée, pour tout $\sigma\in\mathbb Z[n]$, d'un objet $X^\sigma$ de
  $C$.
\item \VerdierSemanticAnchor{ch1-2-1-2-differentials}%
  La donnée, pour tout $\sigma\in\mathbb Z[n]$ et tout $i\in[n]$, d'un
  morphisme :
  \VerdierSemanticAnchor{ch1-2-1-2-differential-map}%
  \begin{equation*}
    d^{\sigma,i}\colon X^\sigma\longrightarrow X^{\sigma+e_i},
  \end{equation*}
  ces morphismes étant soumis aux conditions suivantes :
\end{enumerate}

\VerdierSemanticAnchor{ch1-2-1-2-eq1}
\begin{equation}
  \begin{aligned}
    d^{\sigma+e_i,i}d^{\sigma,i}&=0
      &&\forall\sigma\in\mathbb Z[n],\ \forall i\in[n],\\
    d^{\sigma+e_j,i}d^{\sigma,j}
      &=d^{\sigma+e_i,j}d^{\sigma,i}
      &&\forall\sigma\in\mathbb Z[n],\ \forall i,j\in[n].
  \end{aligned}
  \tag{2.1.2.1}
\end{equation}

\VerdierSemanticAnchor{ch1-2-1-3}
\noindent\textbf{2.1.3.}\quad
Un complexe $n$-uple d'objets de $C$ sera appelé le plus souvent un
\emph{complexe $n$-uple de $C$}. Soit $(X^\sigma,d^{\sigma,i})$ un
complexe $n$-uple de $C$.
\VerdierSemanticAnchor{ch1-2-1-3-component}
L'objet $X^\sigma$ est appelé le \emph{composant de degré $\sigma$} du
complexe $(X^\sigma,d^{\sigma,i})$.
\VerdierSemanticAnchor{ch1-2-1-3-differential}
Le morphisme $d^{\sigma,i}$ est appelé la \emph{différentielle de degré
$\sigma$ dans la direction $i$} du complexe
$(X^\sigma,d^{\sigma,i})$.
\VerdierSemanticAnchor{ch1-2-1-3-notation}
Soit $Y$ un complexe $n$-uple de $C$. Le composant de degré $\sigma$ de
$Y$ est noté $Y^\sigma$. La différentielle de degré $\sigma$ dans la
direction $i$ de $Y$ est notée $d_Y^{\sigma,i}$.

\VerdierSemanticAnchor{ch1-2-1-4}
\VerdierSemanticAnchor{ch1-2-1-4-definition}
\noindent\textbf{Définition 2.1.4.}\quad
Soient $Y$ et $Y'$ deux complexes $n$-uples de $C$. Un \emph{morphisme
de complexes de $Y$ dans $Y'$} est une famille de morphismes de $C$ :

\VerdierSemanticAnchor{ch1-2-1-4-family}
\begin{equation*}
  \bigl(f^\sigma\colon Y^\sigma\longrightarrow Y'^{\sigma}\bigr)_{
    \sigma\in\mathbb Z[n]}
\end{equation*}
qui commutent aux différentielles, \emph{i.e.} telle que :

\VerdierSemanticAnchor{ch1-2-1-4-eq1}
\begin{equation}
  f^{\sigma+e_i}d_Y^{\sigma,i}
  =d_{Y'}^{\sigma,i}f^\sigma.
  \tag{2.1.4.1}
\end{equation}
